\documentclass[11pt]{article}
% Standalone preamble for web-only supplementary notes.
\usepackage[T1]{fontenc}
\usepackage{amsmath,amssymb,amsthm}
\usepackage{hyperref}
\usepackage{cleveref}
\DeclareMathOperator{\Map}{Map}
\DeclareMathOperator{\Fun}{Fun}
\newcommand{\iso}{\xrightarrow{\sim}}
\newtheorem{lemma}{Supplementary Lemma}
\newtheorem{corollary}[lemma]{Supplementary Corollary}
\crefname{lemma}{Supplementary Lemma}{Supplementary Lemmas}
\Crefname{lemma}{Supplementary Lemma}{Supplementary Lemmas}
\crefname{corollary}{Supplementary Corollary}{Supplementary Corollaries}
\Crefname{corollary}{Supplementary Corollary}{Supplementary Corollaries}

\begin{document}

\section*{From a universal law to a particular calculation}

In \href{coherences.html\#post:Composition_Of_Identifications}{Section 1.2},
the laws for composition of identifications are stated as identifications of
functors on animae of inputs. Applying such a law involves a small calculation:
substitute the desired inputs, and compare the resulting endpoints with the
expressions in the statement one wants to prove. This note makes that procedure
explicit.

We use the primitives, terminal category, products, and composition structure
of Sections 1.1--1.2, together with the vertical coherence laws where indicated.
All categories and functors below are absolute. A family with parameter $X$
means an ordinary functor from $X$ into a fixed identification anima.
It does not mean a category varying over $X$.

We write $\beta\cdot\alpha$ for vertical composition, with $\alpha$ applied
first, and $u\star\alpha$ and $\alpha\star r$ for postwhiskering and
prewhiskering. For functor composition we retain $\circ$.

\section*{Lifting two projection comparisons}

The \href{basic-vocabulary.html\#post:Product_Of_Categories}{product axiom}
specifies an equivalence
\[
 q_{F,G}\colon(F=G)\iso
   (\pi_1F=\pi_1G)\times(\pi_2F=\pi_2G)
\]
for functors $F,G\colon X\to C\times D$. Its components are postwhiskering
with the two projections.

\begin{lemma}
\label[lemma]{lem:Supplementary_Product_Lift}
For specified identifications $a\colon\pi_1F=\pi_1G$ and
$b\colon\pi_2F=\pi_2G$, there is an identification $\ell\colon F=G$
equipped with identifications
\[
 \pi_1\star\ell=a,\qquad \pi_2\star\ell=b.
\]
\end{lemma}
\begin{proof}[Proof idea]
Choose the inverse $q_{F,G}^{-1}$ supplied with the equivalence and put
$\ell=q_{F,G}^{-1}\circ\langle a,b\rangle$.
The inverse comparison, the associator, and the left unitor give
\[
 q_{F,G}\circ\ell=\langle a,b\rangle.
\]
Postwhisker this comparison with each projection. The product projection
comparisons, including those defining the components of $q_{F,G}$, give the
two asserted identifications. Thus the construction supplies the projection
computations as well as the lift itself.
\end{proof}

\section*{Substitution through pairing and composition}

For $a\colon X\to C$, $b\colon X\to D$, and $r\colon R\to X$, the
projection comparisons and associator identify the projections of
$\langle a,b\rangle\circ r$ with those of
$\langle a\circ r,b\circ r\rangle$.

\begin{lemma}
\label[lemma]{lem:Supplementary_Substitution_Comparisons}
There are identifications
\[
 \langle a,b\rangle\circ r=\langle a\circ r,b\circ r\rangle.
\]
For families $\sigma\colon X\to(f=g)$ and $\tau\colon X\to(g=h)$,
there is also an identification
\[
 (\tau\cdot\sigma)\circ r
   =(\tau\circ r)\cdot(\sigma\circ r).
\]
More generally, specified comparisons
$\tau\circ r=\tau'$ and $\sigma\circ r=\sigma'$ give
\[
 (\tau\cdot\sigma)\circ r=\tau'\cdot\sigma'.
\]
\end{lemma}
\begin{proof}[Proof idea]
Apply \Cref{lem:Supplementary_Product_Lift} to the two projection comparisons
to obtain the pairing formula. Vertical composition of families is defined
by composing their pairing with the composition functor
$(g=h)\times(f=g)\to(f=h)$. Reassociate, apply the pairing formula, and
postwhisker it with this composition functor. For the last assertion, pair
the two specified comparisons and apply the same composition functor.
\end{proof}

\section*{Keeping the endpoint comparisons}

\begin{lemma}
\label[lemma]{lem:Supplementary_Normalized_Specialization}
Suppose $F,G\colon X\to Y$ are equipped with an identification
$\kappa\colon F=G$. Let $r\colon R\to X$, and suppose that
$F',G'\colon R\to Y$ are equipped with specified identifications
\[
 l\colon F\circ r=F',\qquad t\colon G\circ r=G'.
\]
Then there is an identification $F'=G'$, given by
\[
 \bigl(t\cdot(\kappa\star r)\bigr)\cdot l^{-1}.
\]
\end{lemma}
\begin{proof}[Proof idea]
The inverse of $l$ goes from $F'$ to $F\circ r$, prewhiskering $\kappa$
goes from $F\circ r$ to $G\circ r$, and $t$ goes from $G\circ r$ to $G'$.
Compose these three identifications with the displayed parenthesization.
\end{proof}

\begin{corollary}
\label[corollary]{cor:Supplementary_Specialized_Vertical_Laws}
For absolute identifications $\sigma\colon f=g$, $\tau\colon g=h$, and
$\upsilon\colon h=k$, the universal vertical coherence laws give
\[
 \mathrm{id}_g\cdot\sigma=\sigma,
 \qquad \sigma\cdot\mathrm{id}_f=\sigma,
\]
\[
 (\upsilon\cdot\tau)\cdot\sigma
   =\upsilon\cdot(\tau\cdot\sigma),
\]
\[
 \sigma^{-1}\cdot\sigma=\mathrm{id}_f,
 \qquad \sigma\cdot\sigma^{-1}=\mathrm{id}_g.
\]
These are identifications between the displayed absolute objects of
identification animae.
\end{corollary}
\begin{proof}[Proof idea]
For a unit or inverse law, precompose the universal comparison with
$\sigma\colon *\to(f=g)$. Use
\Cref{lem:Supplementary_Substitution_Comparisons} to compare the substituted
composite with the displayed composite. Constants specialize to their
specified values using uniqueness of functors to $*$ and the unitors.
For associativity, use the point whose three projections are
$\upsilon,\tau,\sigma$, retaining the product projection comparisons at
each step. Apply \Cref{lem:Supplementary_Normalized_Specialization} to the
resulting comparisons of both endpoints.
\end{proof}

The same procedure works with a functor from any parameter category into
the relevant input anima. Each application retains its own endpoint
comparisons. No choice of a universal normal form for all expressions is
needed for these finite calculations.

\end{document}
